\section{Introduction}
Employers routinely receive far more applications than people can read, and software now decides
which resumes reach a human. Screening software is therefore a consequential system, and regulators
have started to treat it as one: the European AI Act places systems that filter or evaluate job
applications in its high-risk category, with obligations on transparency and human oversight
\cite{euaiact2024}. Two properties follow for any practical matcher. A reviewer must be able to
see \emph{why} a candidate was scored as they were, and the builder must be able to show that
irrelevant attributes such as a name do not change the outcome.

Research on person--job fit has mostly pursued accuracy. Convolutional and recurrent matchers learn
from historical application logs \cite{zhu2018pjfnn,qin2018apjfnn}; fine-tuned sentence encoders use
hundreds of thousands of consultant-labelled pairs \cite{lavi2021consultantbert}; and recent work
adds language-model data augmentation \cite{yu2024confit,yu2025confitv2} or multi-agent prompting
\cite{lo2025multiagent}. These systems share three weaknesses for deployment. Their training data
are private, so results cannot be reproduced. Their output is usually one similarity score, with
no statement of which requirement is met by which part of the resume. And bias is treated as a
future direction or left unmeasured, although audits show that embedding retrievers and language models
rank resumes differently when only the name changes
\cite{wilson2024silicon,armstrong2024silicon,an2024discriminate,wang2024jobfair}.
Removing explicit gender or race indicators is a common response, yet it is known to leave proxy
signal behind \cite{dearteaga2019biasinbios,gonen2019lipstick}.

A further, less discussed problem concerns evaluation. The public resume--job fit benchmark we use
contains 8{,}000 labelled pairs but only 643 distinct resumes and 351 distinct job descriptions. Any
single random or official split therefore lets a model profit from having seen the same resume
before. A scorer that never looks at the job description can already be competitive on such a split.

This paper makes four contributions.
\begin{itemize}
\item \textbf{EviMatch}, an open, deterministic matching pipeline that (i) removes identity cues before
  scoring, (ii) parses a posting into weighted requirements, (iii) links each requirement to its
  best-supporting resume sentence, and (iv) combines lexical, semantic, skill and relational
  experience features in a low-capacity head whose every signal can be inspected.
\item A \textbf{leakage analysis} of the benchmark and a doubly-disjoint evaluation protocol in which
  neither the resume nor the job description of a test pair occurs in training.
\item Empirical findings with confidence intervals: which feature groups carry the signal, how little
  labelled data is needed, and that cited evidence is faithful to the model under a deletion test.
\item A \textbf{counterfactual bias audit} for explicit and proxy cues. The audit exposed a defect in our own
  redactor, which we report and fix; the corrected audit quantifies what redaction does and does not remove.
\end{itemize}

Section~\ref{sec:related} reviews related work, Section~\ref{sec:method} describes the system,
Section~\ref{sec:setup} the experimental protocol, Section~\ref{sec:results} the results, and
Section~\ref{sec:conclusion} concludes with limitations and concrete next steps.
